\section{有限覆盖的构造}
\label{sec:fcover-construction}

Arzelà--Ascoli判据只断言有限覆盖存在。要知道需要多少个代表，就要把证明中的“有限位置、有限取值”换成可以计数的具体构造。本节固定一类具体对象：$[0,1]^d$上满足
\[
\lvert f(x)\rvert\leq M,\qquad \lvert f(x)-f(y)\rvert\leq L\lVert x-y\rVert_\infty\qquad(x,y\in[0,1]^d)
\]
的实值函数全体，记为$\mathcal F^d_{L,M}$，其中$L,M>0$，$\lVert x\rVert_\infty=\max_i\lvert x_i\rvert$。误差用$C([0,1]^d)$上的一致距离衡量。若题设给的是欧氏距离下的Lipschitz常数$L'$，由$\lVert x\rVert_2\leq\sqrt d\lVert x\rVert_\infty$可取$L=\sqrt d\,L'$。

上一节已经说明这一类在一致距离下紧。构造的思路是把总误差$\varepsilon$拆成两份预算：一份给位置离散化，只在网格点上记录函数时，网格点之间的变化要受控；另一份给取值量化，网格点上的实数要换成有限个刻度之一。Lipschitz界控制前者，幅值界决定后者需要多少刻度。

\subsection{网格误差与取值误差}
\label{sec:fcover-grid-and-quantization}

\paragraph{网格与多线性插值}

取正整数$m$，网格$G_m$由坐标都属于$\{0,1/m,\ldots,1\}$的点组成，共$(m+1)^d$个节点。网格把$[0,1]^d$分成$m^d$个边长为$1/m$的小立方体单元。只知道函数在节点上的值时，需要一种方法把这些值填满每个单元，并且填出的函数仍然连续。一维时自然的做法是连折线，高维时把折线逐坐标推广。

设$a$是定义在$G_m$上的实值函数。对单元$Q=\prod_{i=1}^d[j_i/m,(j_i+1)/m]$中的点$x$，记局部坐标$t_i=mx_i-j_i\in[0,1]$。\term{多线性插值}（multilinear interpolation）定义为
\begin{equation}
(Ia)(x)=\sum_{e\in\{0,1\}^d}a\Bigl(\frac{j+e}{m}\Bigr)\prod_{i=1}^d w_{e_i}(t_i),\qquad w_1(t)=t,\quad w_0(t)=1-t .
\label{eq:fcover-multilinear}
\end{equation}
这里$e$遍历单元的$2^d$个顶点，顶点$(j+e)/m$的权重是各坐标方向上一维折线权重的乘积。$d=1$时，式\eqref{eq:fcover-multilinear}就是$(1-t)a(j/m)+t\,a((j+1)/m)$，即连接相邻节点值的折线。

这一定义有三条性质。第一，权重非负，且$\sum_e\prod_iw_{e_i}(t_i)=\prod_i\bigl((1-t_i)+t_i\bigr)=1$，所以$(Ia)(x)$是所在单元顶点值的凸组合。第二，在顶点处只有一个权重等于$1$，所以$Ia$在节点上取值$a$。第三，在两个单元的公共面上，例如$t_i=1$时，所有$e_i=0$的顶点权重为零，只剩公共面上的顶点；邻接单元中同一点的局部坐标为$t_i=0$，留下的也正是这些顶点，权重相同。因此两侧公式在公共面上一致，$Ia$是$[0,1]^d$上的连续函数。

\paragraph{两类误差的分别控制}

给定$f\in\mathcal F^d_{L,M}$，把它在节点上的值记为$f|_{G_m}$。再取刻度间距$2\eta$，令$K=\lceil M/\eta\rceil$，刻度为$c_k=-M+(2k-1)\eta$（$k=1,\ldots,K$）。相邻刻度相距$2\eta$，首末刻度满足$c_1=-M+\eta$、$c_K\geq M-\eta$，所以$[-M,M]$中每个数到最近刻度的距离不超过$\eta$。把每个节点值换成最近的刻度（等距时取较小者），得到量化值$q_f$。这一步称为\term{量化}（quantization），它让节点值只剩$K$种可能。

用$Iq_f$作$f$的代表，由三角不等式，
\[
\lVert f-Iq_f\rVert_\infty\leq\underbrace{\lVert f-I(f|_{G_m})\rVert_\infty}_{\text{网格误差}}+\underbrace{\lVert I(f|_{G_m})-Iq_f\rVert_\infty}_{\text{取值误差}} .
\]
两项分别只与一种预算有关。网格误差：设$x$属于某个单元，该单元的每个顶点$v$都满足$\lVert x-v\rVert_\infty\leq1/m$。由凸组合性质，
\[
\lvert f(x)-I(f|_{G_m})(x)\rvert=\Bigl\lvert\sum_v w_v\bigl(f(x)-f(v)\bigr)\Bigr\rvert\leq\sum_vw_v\,\frac Lm=\frac Lm ,
\]
其中$w_v$是式\eqref{eq:fcover-multilinear}中顶点$v$的权重。取值误差：$I$对节点值是线性的，凸组合给出$\lvert(Ia)(x)-(Ib)(x)\rvert\leq\max_{v\in G_m}\lvert a(v)-b(v)\rvert$，而$f|_{G_m}$与$q_f$在每个节点相差不超过$\eta$，所以这一项不超过$\eta$。

\begin{theorem}[Lipschitz类的网格覆盖]
\label{thm:fcover-grid-cover}
对正整数$m$与$\eta>0$，记$K=\lceil M/\eta\rceil$，则
\[
N_{C([0,1]^d)}\Bigl(\frac Lm+\eta,\ \mathcal F^d_{L,M}\Bigr)\leq K^{(m+1)^d}.
\]
特别地，取$m=\lceil2L/\varepsilon\rceil$、$\eta=\varepsilon/2$，得
\begin{equation}
N_{C([0,1]^d)}(\varepsilon,\mathcal F^d_{L,M})\leq\Bigl\lceil\frac{2M}{\varepsilon}\Bigr\rceil^{(\lceil2L/\varepsilon\rceil+1)^d}.
\label{eq:fcover-grid-bound}
\end{equation}
\end{theorem}

\begin{proof}
所有可能的量化节点值$q:G_m\to\{c_1,\ldots,c_K\}$至多有$K^{(m+1)^d}$种，相应的函数$Iq$组成一个有限集。上面的两类误差估计说明，每个$f\in\mathcal F^d_{L,M}$到$Iq_f$的一致距离不超过$L/m+\eta$。代入$m=\lceil2L/\varepsilon\rceil$与$\eta=\varepsilon/2$，有$L/m\leq\varepsilon/2$，总误差不超过$\varepsilon$。
\end{proof}

\paragraph{代表取在类外}

定理中的代表$Iq$取在环境空间$C([0,1]^d)$中，它们未必属于$\mathcal F^d_{L,M}$：量化可能把相邻节点值推向相反方向，使插值函数的斜率超过$L$。例如$d=1$、$L=0.01$、$M=1$、$\varepsilon=0.1$时，$m=1$、$\eta=0.05$，刻度为$\pm0.05,\pm0.15,\ldots$。函数$f(x)=0.01(x-1/2)$在$0$和$1$处的值$\mp0.005$分别量化为$-0.05$和$0.05$，插值折线的斜率为$0.1$，不在类中。

若要求代表本身属于函数类，可以用式\eqref{eq:fcover-center-change}调整半径：$N_{\mathcal F}(\varepsilon,\mathcal F)\leq N_{C([0,1]^d)}(\varepsilon/2,\mathcal F)$，即以$\varepsilon/2$运行上面的构造，再用类中成员替换中心。代价只是把精度参数减半。

以$d=1$、$L=M=1$、$\varepsilon=0.1$为例。此时$m=20$，共$21$个节点，$\eta=0.05$，刻度数$K=20$，式\eqref{eq:fcover-grid-bound}给出$N_{C[0,1]}(0.1,\mathcal F^1_{1,1})\leq20^{21}$，度量熵（自然对数）不超过$21\ln20\approx62.9$。若要求中心在类内，改用$\varepsilon/2=0.05$，得$m=40$、$K=40$，上界为$40^{41}$，度量熵不超过$41\ln40\approx151.2$。

\subsection{覆盖数的维数依赖}
\label{sec:fcover-dimension-dependence}

\paragraph{上界中的四个量}

对式\eqref{eq:fcover-grid-bound}取对数，
\begin{equation}
\log N_{C([0,1]^d)}(\varepsilon,\mathcal F^d_{L,M})\leq\Bigl(\Bigl\lceil\frac{2L}{\varepsilon}\Bigr\rceil+1\Bigr)^d\log\Bigl\lceil\frac{2M}{\varepsilon}\Bigr\rceil .
\label{eq:fcover-entropy-bound}
\end{equation}
右边第一个因子是节点数，第二个因子是每个节点上一个刻度编号的信息量。维数$d$出现在节点数的指数上：精度减半，节点数约乘以$2^d$。Lipschitz界$L$通过网格的细度进入节点数，因而按$(L/\varepsilon)^d$起作用；幅值界$M$只决定刻度数，只在对数中出现。

把这一结果与有限维参数族比较，可以看出Lipschitz类的“大小”。若函数由$p$个有界参数线性决定，度量熵大致按$p\log(1/\varepsilon)$增长，参数个数是固定的（第~\ref{sec:fcover-finite-representations}节）。式\eqref{eq:fcover-entropy-bound}中起参数个数作用的节点数却随精度增长，精度越高需要记录的位置越多，这反映了Lipschitz类是无限维的。下表列出$L=M=1$、$\varepsilon=0.1$时上界随维数的变化。

\begin{table}[!htb]
  \centering
  \caption{$L=M=1$、$\varepsilon=0.1$时式\eqref{eq:fcover-entropy-bound}的各项。每个节点的刻度数都是$20$，节点数按$21^d$增长，度量熵上界（自然对数）随之按同一比例增长。}
  \label{tab:fcover-dimension}
  \begin{tabular}{@{}cccc@{}}
    \toprule
    $d$ & 节点数$21^d$ & 每节点刻度数 & 度量熵上界\\
    \midrule
    $1$ & $21$ & $20$ & $62.9$\\
    $2$ & $441$ & $20$ & $1321.1$\\
    $3$ & $9261$ & $20$ & $27743.5$\\
    \bottomrule
  \end{tabular}
\end{table}

\paragraph{利用相邻位置的约束}

式\eqref{eq:fcover-entropy-bound}是可以逐步核验的上界，却不是最优的。它把各节点上的刻度当作互相独立的选择，而Lipschitz条件实际上约束了相邻节点：相距$1/m$的两个节点，函数值相差不超过$L/m$。利用这一约束，可以去掉主项中的对数因子。

把网格节点排成一列$v_0,v_1,\ldots,v_{n-1}$（$n=(m+1)^d$），使相邻两项只在一个坐标上相差$1/m$。这样的“蛇形”次序存在：$d=1$时按$0,1/m,\ldots,1$排列；若$d-1$维网格已排好，就逐层处理最后一个坐标，每层沿$d-1$维次序走完，相邻层交替使用正向与反向，于是一层的末点与下一层的首点只在最后一个坐标上相差$1/m$。

仍取$m=\lceil2L/\varepsilon\rceil$、$\eta=\varepsilon/2$，则$L/m\leq\eta$。对$f\in\mathcal F^d_{L,M}$，相邻量化值满足
\[
\lvert q_f(v_{r+1})-q_f(v_r)\rvert\leq\eta+\frac Lm+\eta\leq3\eta .
\]
刻度差是$2\eta$的整数倍，所以它只能是$-2\eta$、$0$或$2\eta$。量化值序列因而由首项（$K$种）和后面$n-1$个“降、平、升”的选择决定，至多有$K\cdot3^{n-1}$种，于是
\begin{equation}
\log N_{C([0,1]^d)}(\varepsilon,\mathcal F^d_{L,M})\leq\log\Bigl\lceil\frac{2M}{\varepsilon}\Bigr\rceil+\Bigl(\Bigl(\Bigl\lceil\frac{2L}{\varepsilon}\Bigr\rceil+1\Bigr)^d-1\Bigr)\log3 .
\label{eq:fcover-entropy-improved}
\end{equation}
前面$d=1$、$L=M=1$、$\varepsilon=0.1$的例子中，上界从$62.9$降到$\ln20+20\ln3\approx25.0$。

\paragraph{打包下界}

上界改进到哪一步才算够，要靠下界判断。下界来自第~\ref{chap:metric-spaces}章的打包思路：若类中有一组成员两两距离大于$2\varepsilon$，任何半径$\varepsilon$的闭球至多含其中一个，覆盖数就不小于这组成员的个数。下面在每个小单元里放一个可正可负的隆起。

设$L\geq6\varepsilon$且$\varepsilon\leq M/3$。令$k=\lfloor L/(3\varepsilon)\rfloor$、$h=1/k$，把$[0,1]^d$分成$k^d$个边长$h$的单元，单元$Q$的中心记为$c_Q$。在每个单元上取隆起$\varphi_Q(x)=\max\{0,Lh/2-L\lVert x-c_Q\rVert_\infty\}$，它支撑在$Q$上，在$Q$的边界上为零，高度为$Lh/2$。对每个符号选择$s\in\{-1,1\}^{k^d}$，令$g_s=\sum_Qs_Q\varphi_Q$。

先核对$g_s\in\mathcal F^d_{L,M}$。由$k\geq L/(3\varepsilon)-1\geq L/(6\varepsilon)$得高度$Lh/2=L/(2k)\leq3\varepsilon\leq M$。每个$\varphi_Q$关于$\lVert\cdot\rVert_\infty$的Lipschitz常数为$L$；对任意$x,y$，线段$[x,y]$与每个单元的交是一段子线段，在每段上$g_s$等于单个$\pm\varphi_Q$，而$g_s$连续，所以沿线段逐段累加得$\lvert g_s(x)-g_s(y)\rvert\leq L\lVert x-y\rVert_\infty$。

再看距离。两个不同的符号选择$s\neq s'$在某个单元$Q$上不同，在中心$c_Q$处$\lvert g_s-g_{s'}\rvert=2\cdot Lh/2=Lh$。由$k\leq L/(3\varepsilon)$得$Lh=L/k\geq3\varepsilon>2\varepsilon$。所以这$2^{k^d}$个函数两两距离大于$2\varepsilon$，即使中心可取在整个$C([0,1]^d)$中，也有
\[
\log N_{C([0,1]^d)}(\varepsilon,\mathcal F^d_{L,M})\geq k^d\log2\geq\Bigl(\frac{L}{6\varepsilon}\Bigr)^d\log2 .
\]

\paragraph{上界与最优阶}

把下界与式\eqref{eq:fcover-entropy-improved}放在一起，固定$L,M,d$而令$\varepsilon\to0$时，度量熵夹在$(L/\varepsilon)^d$的两个常数倍之间，这就是它的增长阶；幅值$M$只影响低阶的对数项。Kolmogorov与Tikhomirov系统研究了这类函数类的度量熵，给出了更精确的常数与更多正则性类的结果\scite{kolmogorov1959epsilon}。

本节的经验是：直接的网格量化总能给出可以逐项核验的上界，参数选择与误差来源都一目了然；但要判断上界是否已达最优阶，还需要利用不同位置之间的约束改进编码，并构造相应的打包集给出下界。
